\chapter{\texorpdfstring{Large-\(N\) Two-Dimensional QCD and the Light-Front Bound-State Equation}{Large-N Two-Dimensional QCD and the Light-Front Bound-State Equation}}
\label{ch:large-n-two-dimensional-qcd-light-front}
\ConventionsLorentzian

\section{Regulated Theory and Light-Front Coordinates}

Two-dimensional QCD is a useful solvable gauge-theory laboratory precisely
because the gauge field has no local propagating polarization but still
imposes a nontrivial nonabelian Gauss law.  The model considered in this
chapter has gauge group \(SU(N_c)\), \(N_f\) Dirac fermions in the fundamental
representation, and the same trace convention as the four-dimensional QCD
chapter:
\begin{equation}
\begin{gathered}
  \operatorname{tr}_{\square}(t^at^b)=\delta^{ab},
  \\
  (t^a)^i{}_j(t^a)^k{}_\ell
  =
  \delta^i{}_\ell\delta^k{}_j
  -\frac1{N_c}\delta^i{}_j\delta^k{}_\ell .
\end{gathered}
\label{eq:qcd2-su-n-completeness}
\end{equation}
The second identity is the fundamental completeness relation in this
normalization; it is the same color algebra used in
Chapter~\ref{ch:volume-ii-19}.  The coupling \(g_2\) has mass dimension one,
and the local action is written in the connection normalization
\begin{equation}
  S
  =
  \int\dd^2x\,
  \left[
    -\frac1{4g_2^2}\operatorname{tr}_{\square}(F_{\mu\nu}F^{\mu\nu})
    -\bar q(\gamma^\mu D_\mu+m)q
  \right],
  \qquad
  D_\mu=\partial_\mu-\ii A_\mu .
  \label{eq:qcd2-action}
\end{equation}
For several flavors and masses, \(m\) is replaced by a diagonal mass matrix.
The large-\(N_c\) coupling coordinate used below is
\begin{equation}
  \lambda_2:=g_2^2N_c .
  \label{eq:qcd2-thooft-coupling}
\end{equation}
It has dimensions of mass squared.  Because the generator normalization is
fixed by \eqref{eq:qcd2-su-n-completeness}, no hidden factor of \(1/2\) is
allowed in this definition.

Use light-front coordinates
\begin{equation}
  x^\pm=\frac{x^0\pm x^1}{\sqrt2},
  \qquad
  \dd s^2=-2\,\dd x^+\dd x^- ,
  \label{eq:qcd2-light-front-coordinates}
\end{equation}
and take \(x^+\) as time.  The momentum \(p^+\) is conjugate to \(x^-\), and
a massive on-shell particle has
\[
  2p^+p^-=m^2,\qquad p^+>0 .
\]
The light-front Hamiltonian is \(P^-\), while the invariant mass of a state
with no transverse momentum in two dimensions is
\begin{equation}
  M^2=2P^+P^- .
  \label{eq:qcd2-light-front-mass}
\end{equation}

The manipulations below are made at finite regulator before continuum limits
are discussed.  A convenient regulator is to put \(x^-\) on a circle of
length \(2\pi R\), remove or separately constrain the zero mode, impose a
finite harmonic resolution \(K=RP^+\), and only then study \(K\to\infty\).
Equivalently one may use smooth cutoffs in \(p^+\) and a principal-value
prescription for the instantaneous kernel.  The zero mode of the gauge
constraint is not a harmless detail: it fixes the color-neutral sector and
the global electric-flux data.  The bound-state equation derived below is the
continuum form of the regulated color-singlet two-body problem after this
constraint has been imposed.

\section{Light-Front Gauge Constraint}

Choose light-front gauge
\[
  A_-=0 .
\]
Then \(F_{+-}=-\partial_-A_+\).  The \(A_+\) field is nondynamical; it is the
Lagrange-multiplier remnant of Gauss law written in a gauge where the
constraint can be solved explicitly.  The part of the regulated light-front
action depending on \(A_+\) has the form
\begin{equation}
  S[A_+]
  =
  \int\dd x^+\dd x^-\,
  \left[
    \frac1{2g_2^2}(\partial_-A_+^a)(\partial_-A_+^a)
    +A_+^aJ^{+a}
  \right],
  \label{eq:qcd2-a-plus-action}
\end{equation}
where \(J^{+a}\) is the color current of the dynamical light-front fermion
component.  The sign convention in \eqref{eq:qcd2-a-plus-action} is fixed by
\(\dd s^2=-2\dd x^+\dd x^-\) and by the action
\eqref{eq:qcd2-action}.  Variation of \(A_+^a\) gives
\begin{equation}
  -\frac1{g_2^2}\partial_-^2A_+^a+J^{+a}=0 .
  \label{eq:qcd2-gauss-light-front}
\end{equation}
Let
\[
  G:=(-\partial_-^2)^{-1}
\]
on the zero-mode-free subspace selected by the regulator.  On the infinite
line, \(G\) has kernel \(-|x^-|/2\) as an inverse of \(-\partial_-^2\) on
test functions of total integral zero; on the circle the corresponding
kernel is the periodic zero-average Green function.  Solving
\eqref{eq:qcd2-gauss-light-front} and substituting back gives the positive
instantaneous Coulomb Hamiltonian
\begin{equation}
  P^-_{\rm inst}
  =
  \frac{g_2^2}{2}
  \int\dd x^-\dd y^-\,
  J^{+a}(x^-)\,G(x^--y^-)\,J^{+a}(y^-),
  \label{eq:qcd2-instantaneous-hamiltonian}
\end{equation}
with the zero-mode sector understood as above.

\begin{proposition}[Finite-regulator Gauss-law reduction]
\label{prop:qcd2-gauss-law-reduction}
At finite light-front regulator, after the \(A_+\) zero mode has been either
fixed by the color-neutral Gauss-law constraint or kept as a separate compact
degree of freedom, the nonzero-mode integration over \(A_+\) is an ordinary
Gaussian integration and is equivalent to the instantaneous Hamiltonian
\eqref{eq:qcd2-instantaneous-hamiltonian}.
\end{proposition}

\begin{proof}
On the nonzero-mode regulated space, \(-\partial_-^2\) is a positive
invertible symmetric operator.  Integrating \eqref{eq:qcd2-a-plus-action} by
parts in \(x^-\) gives
\[
  S[A_+]
  =
  \int\dd x^+\,
  \left[
    -\frac1{2g_2^2}\langle A_+,\partial_-^2A_+\rangle
    +\langle A_+,J^+\rangle
  \right].
\]
Completing the square with
\[
  A_+=-g_2^2(-\partial_-^2)^{-1}J^+ + A_+'
\]
separates a Gaussian in \(A_+'\) from the current-current term.  In the
Hamiltonian \(P^-\), the latter appears with the positive kernel \(G\) as in
\eqref{eq:qcd2-instantaneous-hamiltonian}.  The completion of the square does
not determine the zero mode; that mode is exactly the global Gauss-law
constraint and must be part of the regulator datum.
\end{proof}

\section{Planar Color-Singlet Mesons}

Consider one quark flavor of mass \(m_1\) and one antiquark flavor of mass
\(m_2\).  A color-singlet light-front meson state of total momentum \(P^+>0\)
has the leading large-\(N_c\) form
\begin{equation}
  \ket{P^+;\phi}
  =
  \frac1{\sqrt{N_c}}
  \sum_{i=1}^{N_c}
  \int_0^1\dd x\,
  \phi(x)\,
  b_i^\dagger(xP^+)d^{\dagger i}((1-x)P^+)\ket{\vac}
  \label{eq:qcd2-meson-state}
\end{equation}
after a choice of light-front creation-operator normalization.  The function
\(\phi\) is a wavefunction on \(0<x<1\).  Its endpoint behavior is part of
the domain of the mass operator.  For positive renormalized endpoint masses,
normalizability and finiteness of the mass expectation force sufficient
vanishing at \(x=0,1\); in the chiral limit an endpoint-constant wavefunction
is allowed by the subtracted Coulomb kernel below.

The free light-front Hamiltonian contributes
\[
  P^-_{\rm kin}
  =
  \frac1{2P^+}
  \left[
    \frac{m_1^2}{x}+\frac{m_2^2}{1-x}
  \right]
\]
on the component with momentum fraction \(x\).  The instantaneous interaction
\eqref{eq:qcd2-instantaneous-hamiltonian} transfers longitudinal momentum
between the quark and antiquark.  The planar color contraction follows from
\eqref{eq:qcd2-su-n-completeness}: the first term preserves the singlet with
coefficient \(N_c\), while the second term is suppressed by \(N_c^{-1}\) in
the normalized singlet matrix element.  The real exchange kernel and the
quark and antiquark self-energy terms combine into a subtracted principal
value operator.

\begin{definition}[Subtracted 't Hooft operator]
\label{def:qcd2-subtracted-thooft-operator}
For \(\gamma_2>0\) and endpoint masses \(m_1^2,m_2^2\geq0\), define
\begin{equation}
  (\mathcal M^2\phi)(x)
  =
  \left[
    \frac{m_1^2}{x}
    +
    \frac{m_2^2}{1-x}
  \right]\phi(x)
  +
  \gamma_2\,
  \operatorname{PV}\int_0^1
  \frac{\phi(x)-\phi(y)}{(x-y)^2}\,\dd y .
  \label{eq:qcd2-subtracted-thooft-operator}
\end{equation}
The principal value is best defined by the quadratic form in
Proposition~\ref{prop:qcd2-thooft-positive-form}.  For smooth test functions,
it is the finite part obtained by applying the same cutoff to real exchange
and self-energy terms before removing the cutoff.  In the large-\(N_c\) QCD
normalization of this chapter,
\begin{equation}
  \gamma_2=\frac{\lambda_2}{\pi}
  \label{eq:qcd2-thooft-gamma-normalization}
\end{equation}
for the standard continuum principal-value convention.  A finite DLCQ
regulator may instead take \(\gamma_2\) as its displayed matrix coefficient
and use \eqref{eq:qcd2-thooft-gamma-normalization} as the continuum matching
condition.
\end{definition}

\begin{theorem}[Planar light-front meson equation]
\label{thm:qcd2-planar-thooft-equation}
Assume the finite light-front regulator described above, take
\(N_c\to\infty\) with \(\lambda_2=g_2^2N_c\) fixed and \(N_f\) fixed, and
restrict to the color-singlet one-quark-one-antiquark sector with endpoint
masses \(m_1,m_2\) fixed in the chosen renormalization scheme.  If the
regulated planar Hamiltonians converge on this sector to the subtracted
operator of Definition~\ref{def:qcd2-subtracted-thooft-operator}, then the
leading large-\(N_c\) meson masses are the spectral values \(M^2\) satisfying
\begin{equation}
  M^2\phi(x)
  =
  \left[
    \frac{m_1^2}{x}
    +
    \frac{m_2^2}{1-x}
  \right]\phi(x)
  +
  \frac{\lambda_2}{\pi}\,
  \operatorname{PV}\int_0^1
  \frac{\phi(x)-\phi(y)}{(x-y)^2}\,\dd y .
  \label{eq:qcd2-thooft-equation}
\end{equation}
\end{theorem}

\begin{proof}
The free term follows from \(M^2=2P^+P^-\): a quark with momentum \(xP^+\)
contributes \(m_1^2/(2xP^+)\) to \(P^-\), and the antiquark contributes
\(m_2^2/(2(1-x)P^+)\).  Multiplication by \(2P^+\) gives the endpoint mass
operator.

For the interaction, write the current in momentum modes.  At finite
regulator, the kernel \(G=(-\partial_-^2)^{-1}\) contributes a factor
\((k^+)^{-2}\) for a nonzero exchanged longitudinal momentum \(k^+\), with
the zero mode excluded by Gauss law.  Acting on the singlet state
\eqref{eq:qcd2-meson-state}, the color algebra
\eqref{eq:qcd2-su-n-completeness} gives a leading contraction
\(\delta^i{}_\ell\delta^k{}_j\) and a subleading trace subtraction.  Hence the
leading interaction is proportional to \(g_2^2N_c=\lambda_2\).

The real exchange term alone has a nonintegrable \(1/(x-y)^2\) singularity.
The self-energy terms generated by the same instantaneous Hamiltonian are the
terms in which the exchanged longitudinal momentum leaves and returns to the
same fermion line.  With a common regulator for the real and self-energy
pieces, their sum is the finite-part difference
\[
  \operatorname{PV}\int_0^1
  \frac{\phi(x)-\phi(y)}{(x-y)^2}\,\dd y .
\]
The coefficient of this subtracted kernel is \(\lambda_2/\pi\) in the
continuum momentum convention used for light-front creation operators.  The
assumed convergence of the regulated planar Hamiltonians then gives
\eqref{eq:qcd2-thooft-equation}.  The statement is deliberately phrased with
the regulator-convergence hypothesis exposed: the algebra above derives the
finite-regulator reduction and the planar kernel, while the continuum
operator domain is an analytic input.
\end{proof}

\begin{remark}[Relation to the unsubtracted form]
A common finite-part coordinate writes the same regulated planar equation
in the form
\[
  M^2\phi(x)
  =
  \left[
    \frac{\widetilde m_1^2}{x}
    +
    \frac{\widetilde m_2^2}{1-x}
  \right]\phi(x)
  -
  \frac{\lambda_2}{\pi}\,
  \operatorname{PV}\int_0^1
  \frac{\phi(y)}{(x-y)^2}\,\dd y .
\]
This is not a different physical equation; it is a different finite-part
coordinate.  The singular term proportional to \(\phi(x)\) has been assigned
to the endpoint masses \(\widetilde m_i^2\).  The subtracted form
\eqref{eq:qcd2-thooft-equation} displays the real--self-energy pairing and
has a positive quadratic form, so it is the convention used in this
monograph.
\end{remark}

\section{Quadratic Form, Positivity, and the Chiral Limit}

\begin{proposition}[Positive form of the subtracted kernel]
\label{prop:qcd2-thooft-positive-form}
Let \(\phi\) be a real smooth function on \([0,1]\).  Then
\begin{equation}
  \int_0^1\dd x\,\phi(x)
  \operatorname{PV}\int_0^1
  \frac{\phi(x)-\phi(y)}{(x-y)^2}\,\dd y
  =
  \frac12
  \int_0^1\dd x\int_0^1\dd y\,
  \frac{(\phi(x)-\phi(y))^2}{(x-y)^2}.
  \label{eq:qcd2-thooft-positive-form}
\end{equation}
Consequently, for \(m_1^2,m_2^2,\gamma_2\geq0\), the quadratic form of
\(\mathcal M^2\) is nonnegative on its form domain.
\end{proposition}

\begin{proof}
Regulate the singularity by deleting \(|x-y|<\epsilon\).  The regulated
bilinear expression is
\[
  I_\epsilon
  =
  \int_0^1\dd x\int_{|x-y|>\epsilon}\dd y\,
  \phi(x)\frac{\phi(x)-\phi(y)}{(x-y)^2}.
\]
Interchange \(x\) and \(y\) in the same integral.  Since
\((x-y)^2=(y-x)^2\), the average of the two equal expressions is
\[
  I_\epsilon
  =
  \frac12
  \int_{|x-y|>\epsilon}\dd x\,\dd y\,
  \frac{(\phi(x)-\phi(y))^2}{(x-y)^2}.
\]
For smooth \(\phi\), the numerator is \(O((x-y)^2)\) near the diagonal, so
the right hand side has a finite \(\epsilon\downarrow0\) limit.  This proves
\eqref{eq:qcd2-thooft-positive-form}.  Adding the endpoint mass terms gives
nonnegative contributions when \(m_i^2\geq0\).
\end{proof}

In the equal-mass chiral limit \(m_1=m_2=0\), the constant wavefunction
\[
  \phi_0(x)=1
\]
lies in the kernel of the subtracted Coulomb operator.  Therefore
\[
  M_0^2=0
\]
in the leading planar equation.  This massless meson is a special feature of
the large-\(N_c\) two-dimensional model and of the order of limits used here.
It should not be confused with a general theorem about four-dimensional QCD,
nor with the Schwinger-model mass generation of
Chapter~\ref{ch:schwinger-model}; the local gauge dynamics and anomaly
structure are different.

\section{Finite DLCQ Matrix and Numerical Benchmark}

The finite matrix used in
Chapter~\ref{ch:constructive-hamiltonian-truncation-dlcq-benchmarks} is the
DLCQ discretization of the subtracted form above.  Fix \(K\ge3\), grid points
\[
  x_n=\frac nK,\qquad n=1,\ldots,K-1,
\]
and a coefficient \(\gamma_2>0\).  Define
\begin{equation}
  (M_K^2)_{nm}
  =
  \delta_{nm}
  \left[
    \frac{m_1^2}{x_n}
    +
    \frac{m_2^2}{1-x_n}
    +
    \gamma_2K
    \sum_{\substack{j=1\\j\ne n}}^{K-1}\frac1{(n-j)^2}
  \right]
  -
  (1-\delta_{nm})\frac{\gamma_2K}{(n-m)^2}.
  \label{eq:qcd2-thooft-dlcq-matrix}
\end{equation}
Then, for \(v\in\mathbb R^{K-1}\),
\begin{equation}
  v^TM_K^2v
  =
  \sum_{n=1}^{K-1}
  \left[
    \frac{m_1^2}{x_n}
    +
    \frac{m_2^2}{1-x_n}
  \right]v_n^2
  +
  \gamma_2K
  \sum_{1\le n<m\le K-1}
  \frac{(v_n-v_m)^2}{(n-m)^2}.
  \label{eq:qcd2-thooft-dlcq-quadratic-form}
\end{equation}
This identity is the finite-dimensional shadow of
\eqref{eq:qcd2-thooft-positive-form}.  It also fixes the sign of the
off-diagonal matrix elements.  In the massless case, the constant vector is
an exact zero mode of the kernel part at every \(K\), matching the continuum
constant solution.

\begin{controlledapproximation}[DLCQ use of the 't Hooft equation]
A DLCQ calculation of two-dimensional large-\(N_c\) QCD must state the
harmonic resolution \(K\), the treatment of zero modes, the endpoint mass
coordinate, the coefficient \(\gamma_2\) and its matching to
\(\lambda_2/\pi\), the normalization of finite eigenvectors, and the
extrapolation model for \(K\to\infty\).  A finite eigenvalue of
\eqref{eq:qcd2-thooft-dlcq-matrix} is a certified eigenvalue of a regulated
matrix; it becomes a continuum meson mass only after these convergence data
are supplied.
\end{controlledapproximation}

The companion script
\[
  \texttt{calculation-checks/thooft\_model\_checks.py}
\]
verifies the finite color normalization, the quadratic-form identity
\eqref{eq:qcd2-thooft-dlcq-quadratic-form}, positivity for positive endpoint
masses, and the exact massless constant zero mode of the subtracted kernel.

\section{What This Model Does and Does Not Prove}

The model gives an exact large-\(N_c\) bound-state equation in a genuine
confining gauge theory, but its lessons have a precise domain.
\begin{enumerate}
  \item The meson equation follows from a gauge constraint and planar color
  algebra, not from treating quarks as physical external particles.
  \item The absence of transverse gluons in \(1+1\) dimensions is essential.
  Four-dimensional QCD has propagating gluons, running coupling,
  spin-dependent interactions, and pair-creation channels that are not
  captured by \eqref{eq:qcd2-thooft-equation}.
  \item Large-\(N_c\) factorization makes single-meson sectors stable at
  leading order.  Meson interactions and widths enter at subleading powers of
  \(1/N_c\).
  \item The light-front derivation is exact only after the zero-mode and
  endpoint-mass conventions are fixed.  Omitting those conventions changes the
  mathematical problem.
\end{enumerate}
Thus two-dimensional large-\(N_c\) QCD is a benchmark for confinement,
bound-state equations, and light-front regularization.  It is not a substitute
for a nonperturbative construction of four-dimensional QCD, but it is one of
the cleanest places where the relation between gauge constraints,
color-singlet Hilbert-space sectors, and a computable meson spectrum can be
studied without appealing to perturbative \(S\)-matrix intuition.
